\subsection{Retracts, One-Sided Inverses, and Split Maps}

Throughout this subsection, $R$ denotes a nonzero commutative ring with identity.

\begin{theorem} \label{thm:direct-summands-extend-epimorphisms} \hlt{Direct Summands Extend Epimorphisms}\\
Let $N\leq M$ and let $M_1$ be an $R$-module.
\begin{enumerate}[label=(\roman*)]
\setlength{\itemsep}{0pt}
\item If $M=N\oplus H$ and $\sigma:N\twoheadrightarrow M_1$, then $\exists$ a unique $\overline{\sigma}\in\operatorname{Hom}_R(M,M_1)$ such that\\
$\overline{\sigma}(n+h)\defeq\sigma(n)$, $n\in N$, $h\in H$, $\overline{\sigma}|_N=\sigma$, $\overline{\sigma}:M\twoheadrightarrow M_1$, and $\ker\overline{\sigma}=\ker\sigma\oplus H$.
\item If $\sigma:N\xrightarrow{\sim}M_1$, then $\mathcal{E}_{\sigma}\defeq\{\overline{\sigma}:M\twoheadrightarrow M_1:\overline{\sigma}|_N=\sigma\}$ $\longleftrightarrow$ $\mathcal{C}_N(M)\defeq\{H\leq M:M=N\oplus H\}$ via $\overline{\sigma}\mapsto\ker\overline{\sigma}$ and $H\mapsto\overline{\sigma}_H$, where $\overline{\sigma}_H(n+h)\defeq\sigma(n)$.
\end{enumerate}
\begin{proof}
\begin{enumerate}[label=(\roman*)]
\setlength{\itemsep}{0pt}
\item Direct-sum uniqueness makes $\overline{\sigma}$ well-defined, and componentwise computation gives $R$-linearity.\\
Moreover, $\operatorname{im}\overline{\sigma}=\operatorname{im}\sigma=M_1$; $\ker\overline{\sigma}=\{n+h:\sigma(n)=0\}=\ker\sigma\oplus H$. This shows uniqueness.
\item If $H\in\mathcal{C}_N(M)$, then (i) gives $\overline{\sigma}_H\in\mathcal{E}_{\sigma}$ and, since $\ker\sigma=\{0\}$, $\ker\overline{\sigma}_H=H$.\\
Conversely, if $\overline{\sigma}\in\mathcal{E}_{\sigma}$, then $N\cap\ker\overline{\sigma}=\{0\}$, and for each $m\in M$, $n\defeq\sigma^{-1}(\overline{\sigma}(m))\in N$, $m-n\in\ker\overline{\sigma}$.\\
Thus $M=N\oplus\ker\overline{\sigma}$, so the two displayed assignments are inverse.
\end{enumerate}
\hlt{Structural Commentary}: A complement supplies the directions on which an extension may vanish.\\
\hlt{Key Insight}: Extending an isomorphism is equivalent to choosing its kernel complement.
\end{proof}
\end{theorem}

\begin{definition} \label{def:module-retract} \hlt{Retract}\\
Let $N\leq M$.
\begin{enumerate}[label=(\roman*)]
\setlength{\itemsep}{0pt}
\item A \hlt{retraction} onto $N$ is a map $\rho\in\operatorname{Hom}_R(M,N)$ such that $\rho|_N=\operatorname{id}_N$.
\item The submodule $N$ is a \hlt{retract} of $M$ iff $\exists\,\rho\in\operatorname{Hom}_R(M,N)$ with $\rho|_N=\operatorname{id}_N$.
\end{enumerate}
\end{definition}

\begin{corollary} \label{cor:retract-iff-direct-summand-module} \hlt{Retracts Are Direct Summands}\\
For $N\leq M$, $N$ is a retract of $M$ $\Longleftrightarrow$ $\exists\,H\leq M$ such that $M=N\oplus H$.
\begin{proof}
\begin{enumerate}[label=\roman*.]
\setlength{\itemsep}{0pt}
\item[$(\Rightarrow)$] Let $\rho:M\to N$ satisfy $\rho|_N=\operatorname{id}_N$. For every $m\in M$, $m=\rho(m)+(m-\rho(m))$, $\rho(m)\in N$, $m-\rho(m)\in\ker\rho$.\\
Moreover $N\cap\ker\rho=\{0\}$, so $M=N\oplus\ker\rho$.
\item[$(\Leftarrow)$] If $M=N\oplus H$, define $\rho:M\to N$ by $\rho(n+h)\defeq n$, $n\in N$, $h\in H$.\\
Then $\rho\in\operatorname{Hom}_R(M,N)$ and $\rho|_N=\operatorname{id}_N$, so $N$ is a retract.
\end{enumerate}
\hlt{Structural Commentary}: Retractions are projections onto submodules.\\
\hlt{Key Insight}: A submodule is split exactly when it can be recovered by an idempotent map.
\end{proof}
\end{corollary}

\begin{definition} \label{def:one-sided-inverses-module-maps} \hlt{One-Sided Inverses}\\
Let $\tau\in\operatorname{Hom}_R(A,B)$.
\begin{enumerate}[label=(\roman*)]
\setlength{\itemsep}{0pt}
\item A \hlt{left inverse} of $\tau$ is $\lambda\in\operatorname{Hom}_R(B,A)$ such that $\lambda\tau=\operatorname{id}_A$.
\item A \hlt{right inverse} of $\tau$ is $\rho\in\operatorname{Hom}_R(B,A)$ such that $\tau\rho=\operatorname{id}_B$.
\item A \hlt{two-sided inverse} of $\tau$ is $\upsilon\in\operatorname{Hom}_R(B,A)$ such that $\upsilon\tau=\operatorname{id}_A$ and $\tau\upsilon=\operatorname{id}_B$.
\end{enumerate}
\end{definition}

\begin{theorem} \label{thm:module-one-sided-inverses-and-complements} \hlt{One-Sided Inverses and Complements}\\
Let $\sigma:M\to M_1$ be an $R$-homomorphism.\\
For $S\leq N$, set $\mathcal{C}_S(N)\defeq\{H\leq N:N=S\oplus H\}$, $\mathcal{L}_\sigma\defeq\{\lambda\in\operatorname{Hom}_R(M_1,M):\lambda\sigma=\operatorname{id}_M\}$,\\
and $\mathcal{R}_\sigma\defeq\{\rho\in\operatorname{Hom}_R(M_1,M):\sigma\rho=\operatorname{id}_{M_1}\}$.
\begin{enumerate}[label=(\roman*)]
\setlength{\itemsep}{0pt}
\item If $\sigma$ is injective, then $\mathcal{L}_\sigma\xrightarrow{\sim}\mathcal{C}_{\operatorname{im}\sigma}(M_1)$ by $\lambda\mapsto\ker\lambda$.
\item If $\sigma$ is surjective, then $\mathcal{R}_\sigma\xrightarrow{\sim}\mathcal{C}_{\ker\sigma}(M)$ by $\rho\mapsto\operatorname{im}\rho$.
\end{enumerate}
\begin{proof}
\begin{enumerate}[label=(\roman*)]
\setlength{\itemsep}{0pt}
\item For $H\in\mathcal{C}_{\operatorname{im}\sigma}(M_1)$, define $\lambda_H(\sigma m+h)\defeq m$. Directness and injectivity give well-definedness; $\lambda_H\sigma=\operatorname{id}_M$ and $\ker\lambda_H=H$. Conversely, if $\lambda\in\mathcal{L}_\sigma$, then $\operatorname{im}\sigma\cap\ker\lambda=\{0\}$, since $\sigma m\in\ker\lambda\Rightarrow m=0$, and $y=\sigma\lambda y+(y-\sigma\lambda y)$ with $y-\sigma\lambda y\in\ker\lambda$. Thus $M_1=\operatorname{im}\sigma\oplus\ker\lambda$, and the constructions are inverse.
\item For $H\in\mathcal{C}_{\ker\sigma}(M)$, $\sigma|_H:H\to M_1$ is an isomorphism; let $\rho_H=(\sigma|_H)^{-1}$ followed by $H\hookrightarrow M$.\\
Then $\rho_H\in\mathcal{R}_\sigma$ and $\operatorname{im}\rho_H=H$. Conversely, if $\rho\in\mathcal{R}_\sigma$, then $\rho$ is injective, $\operatorname{im}\rho\cap\ker\sigma=\{0\}$, and $m=(m-\rho\sigma m)+\rho\sigma m$ with $m-\rho\sigma m\in\ker\sigma$, $\rho\sigma m\in\operatorname{im}\rho$.\\
Thus $M=\ker\sigma\oplus\operatorname{im}\rho$, and the constructions are inverse.
\end{enumerate}
\hlt{Structural Commentary}: One-sided inverses exist exactly when the relevant image or kernel splits.\\
\hlt{Key Insight}: Injective maps need a complemented image; surjective maps need a complemented kernel.
\end{proof}
\end{theorem}

\begin{proposition} \label{prop:split-epimorphism-modules} \hlt{Split Epimorphisms Split the Domain}\\
Let $\sigma:M\twoheadrightarrow N$ be an $R$-homomorphism and set $\mathcal{R}_\sigma\defeq\{\rho\in\operatorname{Hom}_R(N,M):\sigma\rho=\operatorname{id}_N\}$. These are equivalent:
\begin{enumerate}[label=(\roman*)]
\setlength{\itemsep}{0pt}
\item $\mathcal{R}_\sigma\neq\emptyset$.
\item $\exists\,H\leq M$ such that $M=\ker\sigma\oplus H$.
\end{enumerate}
Moreover, $\rho\in\mathcal{R}_\sigma\Longrightarrow M=\ker\sigma\oplus\operatorname{im}\rho$.
\begin{proof}
\begin{enumerate}[label=\roman*.]
\setlength{\itemsep}{0pt}
\item[$(\Rightarrow)$] Let $\rho\in\mathcal{R}_\sigma$. If $\rho n=0$, then $n=\sigma\rho n=0$, so $\rho$ is injective.\\
If $\rho n\in\ker\sigma$, then $n=\sigma\rho n=0$, hence $\operatorname{im}\rho\cap\ker\sigma=\{0\}$.\\
For $m\in M$, $m=(m-\rho\sigma m)+\rho\sigma m$, with $m-\rho\sigma m\in\ker\sigma$ and $\rho\sigma m\in\operatorname{im}\rho$. Thus $M=\ker\sigma\oplus\operatorname{im}\rho$.
\item[$(\Leftarrow)$] Let $M=\ker\sigma\oplus H$. Then $\ker(\sigma|_H)=H\cap\ker\sigma=\{0\}$, and $\operatorname{im}(\sigma|_H)=N$ since $n=\sigma m$, $m=k+h$, gives $n=\sigma h$. Thus $\sigma|_H:H\xrightarrow{\sim}N$, and $\rho\defeq\iota_H(\sigma|_H)^{-1}\in\mathcal{R}_\sigma$, where $\iota_H:H\hookrightarrow M$.
\end{enumerate}
\hlt{Structural Commentary}: A split epimorphism is a quotient map with complemented kernel.\\
\hlt{Key Insight}: The right inverse chooses canonical representatives for the
cosets of $\ker\sigma$.
\end{proof}
\end{proposition}
